\SourcePage{72}
\SGAchapter{XII}{二次超曲面}{Quadriques}
\begin{center}P. Deligne\end{center}

\textbf{0.0。} 本讲汇集关于\Term{quadric}{二次超曲面}的一些熟知结果；我们在证明 \Term{picard-lefschetz-formula}{Picard--Lefschetz 公式}时将用到这些结果。

\textbf{0.1。} 设 $m$ 是\Term{scheme}{概形} $S$ 上一个 $\mathscr O_S$-模 $M$ 的\Term{section}{截面}。如果对每个 $s\in S$，$m$ 在由\Term{base-change}{基变换} $\{s\}\to S$ 从 $M$ 得到的 $k(s)$-\Term{vector-space}{向量空间} $M_s$ 中的像都非零，就称 $m$ 在 $S$ 上\Term{nowhere-zero}{处处非零}。对 $S=\operatorname{Spec}(A)$ 及一个 $A$-模 $M$ 的元素 $m$，也使用同样的术语。

\SGAsection{1}{二次型}
\textbf{1.1。} 设 $S=\operatorname{Spec}(A)$ 是\Term{affine-scheme}{仿射概形}，$V$ 是一个\Term{rank}{秩}为 $n$ 的\Term{locally-free}{局部自由} $A$-模，$Q$ 是 $V$ 上一个\Term{quadratic-form}{二次型}（Bourb. Alg 9，§3，n\textsuperscript{o}4）。回忆：如果相伴的\Term{bilinear-form}{双线性型}
\[
\Phi(x,y)=Q(x+y)-Q(x)-Q(y)
\]
是\Term{nondegenerate}{非退化}的，即给出 $V$ 与其\Term{dual}{对偶} $V^*$ 之间的一个\Term{isomorphism}{同构}，就称 $Q$ \Term{nondegenerate}{非退化}。如果 $2$ 在 $A$ 中不\Term{invertible}{可逆}，这只有在 $n$ 为偶数时才可能发生（同上）。

\SourcePage{73}
$V$ 上的\Term{quadratic-form}{二次型}构成一个\Term{projective-module}{投射} $A$-模；\Term{quadratic-form}{二次型} $Q$ 在 $S$ 上“\Term{nowhere-zero}{处处非零}”（0.1），当且仅当各个 $Q(v)$ 生成 $A$ 的\Term{ideal}{理想} $A$。

记号 $\mathbb P(V)$ 按 EGA 的意义使用；回忆，若 $A$ 是域，则 $\mathbb P(V)$ 是\Term{rational-point}{有理点}集为\Term{quotient-set}{商集} $(V^*-\{0\})/A^*$ 的\Term{projective-space}{射影空间}。

\Term{quadratic-form}{二次型} $Q$ 可视为 $S$ 上\Term{projective-space}{射影空间} $\mathbb P(V^*)$ 的\Term{invertible-sheaf}{可逆层} $\mathscr O(2)$ 的一个\Term{section}{截面}。当 $Q$ \Term{nowhere-zero}{处处非零}时，（“\Term{homogeneous}{齐次}”）方程 $Q=0$ 定义 $\mathbb P(V^*)$ 的一个\Term{subscheme}{子概形}，它在 $S$ 上\Term{flat}{平坦}，且纯\Term{relative-dimension}{相对维数}为 $n-2$。这个\Term{scheme}{概形}称为 $Q$ 定义的\Term{quadric}{二次超曲面}。若\Term{quadratic-form}{二次型} $Q$ 在 $S$ 的任一点都不为零，且它定义的\Term{quadric}{二次超曲面}在 $S$ 上\Term{smooth}{光滑}，就称 $Q$ 是\Term{ordinary-quadratic-form}{普通的}（ordinaire）。

$Q$ 为\Term{ordinary-quadratic-form}{普通二次型}的充要条件是：从 $A$ 到任意一个域作\Term{scalar-extension}{标量扩张}所得的\Term{quadratic-form}{二次型}都是\Term{ordinary-quadratic-form}{普通的}。若 $A$ 为域，则：

\noindent a) 当 $n$ 为偶数或 $\operatorname{car}(A)\ne2$ 时：$Q$ \Term{ordinary-quadratic-form}{普通} $\Longleftrightarrow Q$ \Term{nondegenerate}{非退化}；

\noindent b) 当 $n$ 为奇数且 $\operatorname{car}(A)=2$ 时：$Q$ \Term{ordinary-quadratic-form}{普通}，当且仅当相伴的（\Term{alternating-form}{交错}）\Term{bilinear-form}{双线性型} $\Phi$ 的\Term{kernel}{核} $N$ 的\Term{dimension}{维数}为一，而且 $Q$ 在 $N$ 上不为零。

\begin{SGAresult}{命题}{1.2}
若 $Q$ 是\Term{ordinary-quadratic-form}{普通的}，且 $n=2m$（相应地，$n=2m+1$），则在 $S$ 上按\Term{etale-topology}{平展拓扑}\Term{local-scheme}{局部}地，$V$ 有一个基 $e$，使得
\SourcePage{74}
\[
Q\left(\sum_{1}^{n}x_ie_i\right)=\sum_{1}^{m-1}x_ix_{i+m}
\]
（相应地，
\[
Q\left(\sum_{1}^{n}x_ie_i\right)=\sum_{1}^{m-1}x_ix_{i+m}+\lambda x_{2m+1}^{2}\qquad(\lambda\text{ 可逆})\text{）。}
\]
\end{SGAresult}
\begin{proof}
对 $m$ 作归纳来证明 1.2。若 $m=0$，断言显然。若 $m>0$，方程 $Q=0$ 的\Term{quadric}{二次超曲面}在 $S$ 上\Term{smooth}{光滑}，且\Term{geometric-fibre}{几何纤维}\Term{nonempty}{非空}；因此它按\Term{etale-topology}{平展拓扑}\Term{local-scheme}{局部}地有\Term{section}{截面}。\Term{local-scheme}{局部}地，这样一个\Term{section}{截面}由一个元素 $e\in V$ 定义，它在 $\operatorname{Spec}(A)$ 的任一点都不为零，且 $Q(e)=0$。在 $S$ 的任一点，$e$ 都不在 $\Phi$ 的\Term{kernel}{核}中；因而\Term{local-scheme}{局部}地可找到 $f'$ 使 $\Phi(e,f')=1$。若 $f=-Q(f')e+f$，则
\[
Q(e)=Q(f)=0,\qquad\Phi(e,f)=1.
\]
若 $V_2$ 是由 $e$ 和 $f$ 生成的\Term{submodule}{子模} $V_1$ 的\Term{orthogonal-complement}{正交补}，则 $V=V_1\oplus V_2$；将归纳假设用于 $V_2$，即得结论。
\end{proof}

\textbf{1.3。} 关于 $Q$ 的\Term{clifford-algebra}{Clifford 代数}的定义及主要性质，参见 Bourb. Alg 9，§9。以 $\rho$ 表示从 $V$ 到 $C(Q)$ 的\Term{canonical-map}{典范映射}。

\begin{SGAresult}{\Term{lemma}{引理}}{1.3.1}
在 1.1 的假设下，对 $\lambda\in A$ 存在唯一的\Term{homomorphism}{同态}
\[
[\lambda]:C^+(\lambda Q)\longrightarrow C^+(Q)
\]
使得对 $x,y\in V$ 有
\[
[\lambda](\rho(x)\rho(y))=\lambda\rho(x)\rho(y).
\]
\end{SGAresult}
\SourcePage{75}
\begin{proof}
这一\Term{homomorphism}{同态}是由 $E$ 的\Term{tensor-algebra}{张量代数}的部分 $T^+(E)$ 上的\Term{endomorphism}{自同态} $[\lambda]$ 通过取商得到的；它在 $T^{2n}(E)$ 上等于 $\lambda^n$。事实上，
\[
[\lambda](x\otimes x-\lambda Q(x))=\lambda(x\otimes x-Q(x)).
\]
\end{proof}

假设 $Q$ 在 $S$ 上\Term{nowhere-zero}{处处非零}。于是\Term{lemma}{引理} 1.2.1 对各个 $\lambda\in A^*$ 定义了\Term{algebra}{代数} $C^+(\lambda Q)$ 之间的一个传递的\Term{isomorphism}{同构}系。

因此，\Term{algebra}{代数} $C^+(Q)$ 仅依赖于 $V$ 以及相差一个因子的 $Q$（在唯一\Term{isomorphism}{同构}意义下）；换言之，$C^+(Q)$ 仅依赖于 $V$ 和 $\mathbb P(V^*)$ 中方程 $Q=0$ 定义的\Term{subscheme}{子概形}。每一个比率\Term{invertible}{可逆}的\Term{homothety}{位似}都是这对\Term{object}{对象} $(V,X\subset\mathbb P(V^*))$ 的一个\Term{automorphism}{自同构}；可以验证，这些\Term{automorphism}{自同构}在 $C^+(V,X)$ 上的\Term{action}{作用}平凡。因此，\Term{algebra}{代数} $C^+(V,X)$ 仅依赖于 $A$ 上的一对\Term{scheme}{概形} $X\subset\mathbb P(V^*)$。我们将它记作 $C^+(X,\mathbb P(V^*))$。

\textbf{1.4。} 设 $V$ 是\Term{rank}{秩}为 $2m>0$ 的\Term{locally-free}{局部自由} $A$-模，$Q$ 是 $V$ 上一个\Term{nondegenerate}{非退化}\Term{quadratic-form}{二次型}。给定 $V$ \Term{resolution}{分解}为两个\Term{totally-singular}{全奇异}（totalement singuliers）\Term{vector-subspace}{子空间}的\Term{resolution}{分解} $V=W_1\oplus W_2$。于是型 $Q(w_1+w_2)$ 给出 $W_1$ 与 $W_2$ 的\Term{perfect-duality}{完美对偶}。回忆：对 $f\in W_2$（将 $W_2$ 与 $W_1^*$ 等同），与 $f$ 的\Term{contraction}{收缩积}是 $\bigwedge W_1$ 上由下式定义的\Term{degree}{次数}为 $-1$ 的\Term{endomorphism}{自同态}：
\[
f\mathbin{\lrcorner}(e_1\wedge\cdots\wedge e_p)
=\sum(-1)^i\langle e_i,f\rangle e_0\wedge\cdots\wedge\widehat e_i\wedge\cdots\wedge e_p.
\]
在 $A$-模 $\bigwedge W_1$ 上赋予 $\mathbb Z/2$-\Term{grading}{分次}，使得
\SourcePage{76}
\[
(\bigwedge W_1)^+=\bigoplus_i\bigwedge^{2i}W_1,
\qquad(\bigwedge W_1)^-=\bigoplus_i\bigwedge^{2i+1}W_1.
\]
于是\Term{algebra}{代数} $\operatorname{End}(\bigwedge W_1)$ 也有 $\mathbb Z/2$-\Term{grading}{分次}：偶（相应地，奇）\Term{endomorphism}{自同态}是保持\Term{grading}{分次}（相应地，\Term{degree}{次数}为 $1\pmod2$）的\Term{endomorphism}{自同态}。

Bourb. Alg 9，§9，n\textsuperscript{o}4 证明了存在唯一的 $\mathbb Z/2$-\Term{graded-algebra}{分次代数}\Term{isomorphism}{同构} $s:C(Q)\xrightarrow{\sim}\operatorname{End}(\bigwedge W_1)$，使得
\[
s(\rho((e,0)))=e\wedge,\qquad s(\rho((0,f)))=f\mathbin{\lrcorner}.
\]
暂以 $e(W_1,W_2)$ 表示 $C^+(Q)$ 中这样的元素：$s(e(W_1,W_2))$ 在 $(\bigwedge W_1)^+$ 上为\Term{identity-map}{恒等映射}，在 $(\bigwedge W_1)^-$ 上为零。$C^+(Q)$ 的\Term{centre-algebra}{中心} $Z(Q)$ \Term{isomorphism}{同构}于 $A\oplus A$，这一\Term{resolution}{分解}由\Term{idempotent}{幂等元} $e(W_1,W_2)$ 和 $1-e(W_2,W_1)$ 定义。此外，$C^+(Q)$ 是 $Z(Q)$ 上的一个\Term{matrix-algebra}{矩阵代数}。

由此以及 1.2，与上述文献中一样可得：

\begin{SGAresult}{命题}{1.5}
在 1.1 的假设下，若 $n=2m>0$ 且 $Q$ \Term{nondegenerate}{非退化}，则 $C^+(Q)$ 的\Term{centre-algebra}{中心} $Z(Q)$ 是 $A$ 上一个\Term{rank}{秩}为 $2$ 的\Term{locally-free}{局部自由}\Term{etale-algebra}{平展代数}，且 $C^+(Q)$ 是 $Z(Q)$ 上的\Term{azumaya-algebra}{Azumaya 代数}。
\end{SGAresult}

\textbf{1.6。} 设 $V$ 和 $Q$ 如 1.4，$W$ 是 $V$ 中一个\Term{rank}{秩}为 $m$ 的\Term{locally-free}{局部自由}\Term{totally-singular}{全奇异}\Term{direct-summand}{直和因子}。设 $W'_1$ 是 $W$ 的一个\Term{complementary-subspace}{补空间}：$V=W\oplus W'_1$。$W$ 的任一个\Term{complementary-subspace}{补空间} $W'_2$ 都可视为一个映射 $u:W'_1\to W$ 的\Term{graph}{图}。型 $\Phi$ 使 $W$ 与 $W'_1$ \Term{dual}{对偶}；因而映射 $u$ 完全由 $W'_1$ 上的\Term{bilinear-form}{双线性型}
\SourcePage{77}
\[
B(x,y)=\Phi(x,u(y))
\]
确定；而 $W'_2$ 是\Term{totally-singular}{全奇异}的，当且仅当在 $W'_1$ 上有
\[
B(x,x)=-Q(x).
\]
由于这样的型 $B$ 总是存在，$W$ 有\Term{totally-singular}{全奇异}的\Term{complementary-subspace}{补空间}。选定其中一个，记为 $W'_0$，则 $W$ 的\Term{totally-singular}{全奇异}\Term{complementary-subspace}{补空间} $W'$ 的\Term{set}{集合}可与 $W'_0$ 上\Term{alternating-form}{交错型}的\Term{set}{集合}等同，也就是与 $S$ 上\Term{affine-space}{仿射空间} $\mathbf E(\bigwedge^2W')$ 的\Term{section}{截面}集等同。对 $\mathbf E(\bigwedge^2W')$ 的每个\Term{section}{截面} $x$，以 $W'_0(x)$ 表示 $W$ 的相应\Term{complementary-subspace}{补空间}，以 $e(x)$ 表示\Term{idempotent}{幂等元} $e(W,W'_0(x))$。函数 $x\mapsto e(x)$ 与任意\Term{scalar-extension}{标量扩张}相容，因而定义了一个从 $\mathbf E(\bigwedge^2W')$ 到 $\operatorname{Spec}(Z(Q))$ 的 $S$-\Term{scheme}{概形}\Term{morphism}{态射} $e$。$S$-\Term{scheme}{概形} $\mathbf E(\bigwedge^2W')$ \Term{flat}{平坦}且\Term{geometric-fibre}{几何纤维}\Term{connected}{连通}，而 $\operatorname{Spec}(Z(Q))$ 在 $S$ 上\Term{etale}{平展}。因此\Term{morphism}{态射} $e$ 经由 $\mathbf E(\bigwedge^2W')$ 在 $S$ 上的一个\Term{section}{截面}\Term{resolution}{分解}，即 $e(x)$ 与 $x$ 无关。这说明了下列记号的合理性。

\begin{SGAresult}{记号}{1.7}
设 $W$ 如上，$W'$ 是 $W$ 的任一个\Term{totally-isotropic}{全迷向}（totalement isotrope）\Term{complementary-subspace}{补空间}。以 $e(W)$ 表示\Term{idempotent}{幂等元} $e(W,W')\in Z(Q)$。
\end{SGAresult}

\textbf{1.8。} 在 1.6 中已经看到：对\Term{invertible}{可逆}的 $\lambda$，$C^+(Q)$ 与 $C^+(\lambda Q)$ 等同。与上面类似的\Term{continuity}{连续性}论证表明，这一等同将
\[
e(W)\in Z(Q)\quad\text{变为}\quad e(W)\in Z(\lambda Q).
\]

\SourcePage{78}
\textbf{1.9。} 回忆，若 $B$ 和 $C$ 是两个 $\mathbb Z/2$-\Term{graded-algebra}{分次代数}，则它们的\Term{graded-tensor-product}{分次张量积} $B\otimes' C$ 中的\Term{multiplication}{乘法}由 Koszul 规则定义：当 $b_i$ 和 $c_i$ 都是\Term{homogeneous}{齐次}的时，
\[
(b_1\otimes c_1)(b_2\otimes c_2)
=(-1)^{\deg(c_1)\deg(b_2)}b_1b_2\otimes c_1c_2.
\]
若 $(V_1,Q_1)$ 和 $(V_2,Q_2)$ 是两个\Term{quadratic-module}{二次模}，则
\[
C(Q_1\oplus Q_2)=C(Q_1)\otimes' C(Q_2)
\]
（Bourb.，同上）。

若 $P$ 和 $Q$ 是两个\Term{locally-free}{局部自由}、在 $A$ 上\Term{finite-type}{有限型}的 $\mathbb Z/2$-\Term{grading}{分次} $A$-模，则按下述规则将 $\mathbb Z/2$-\Term{graded-algebra}{分次代数} $\operatorname{End}(P)\otimes'\operatorname{End}(Q)$ 和 $\operatorname{End}(P\otimes Q)$ 等同：对\Term{homogeneous}{齐次}的 $u,v,p,q$，令
\[
(u\otimes v)(p\otimes q)=(-1)^{\deg(v)\deg(p)}u(p)\otimes v(q).
\]

\begin{SGAresult}{\Term{lemma}{引理}}{1.10}
设 $(V,Q)$ 是两个\Term{quadratic-module}{二次模} $(V_1,Q_1)$、$(V_2,Q_2)$ 的和，其中 $V_i$ \Term{locally-free}{局部自由}且\Term{rank}{秩}为 $2m_i$，$Q_i$ \Term{nondegenerate}{非退化}。设 $V_i=W_i\oplus W'_i$ 是 $V_i$ \Term{resolution}{分解}为两个\Term{totally-singular}{全奇异}\Term{vector-subspace}{子空间}的\Term{resolution}{分解}（$i=1,2$）。则下图交换：
\[
\begin{tikzcd}[column sep=small]
C(Q)\arrow[rr,"\sim"]\arrow[d,"\sim"']&&C(Q_1)\otimes'C(Q_2)\arrow[d,equal]\\
\operatorname{End}(\bigwedge(W_1\oplus W_2))\arrow[r,"\sim"]&
\operatorname{End}(\bigwedge W_1\otimes\bigwedge W_2)\arrow[r,"\sim"]&
\operatorname{End}(\bigwedge W_1)\otimes'\operatorname{End}(\bigwedge W_2).
\end{tikzcd}
\]
\end{SGAresult}
\SourcePage{79}
验证留给读者。由此\Term{lemma}{引理}立即得到
\[
e(W_1\oplus W_2)=e(W_1)e(W_2)+(1-e(W_1))(1-e(W_2)).\tag{1.10.1}
\]

\textbf{1.11。} 下述构造取自 Micali 和 Villamayor。若 $p:Z\to S$ 是 $S$ 的一个双重\Term{etale-covering}{平展覆盖}，则可以把 $Z$ 看作 $S$ 上的一个 $\mathbb Z/2$-\Term{torsor}{挠子}。以 $*$ 表示\Term{torsor}{挠子}的加法；对 $Z$ 的一个\Term{section}{截面} $s$，以 $e(s)$ 表示 $p_*\mathscr O_Z$ 中在 $s$ 上等于 $1$、在别处等于 $0$ 的\Term{idempotent}{幂等元}。沿用 1.10 的记号，存在唯一的\Term{isomorphism}{同构}
\[
\operatorname{Spec}(Z(Q_1))*\operatorname{Spec}(Z(Q_2))\longrightarrow\operatorname{Spec}(Z(Q)),\tag{1.11.1}
\]
\[
+:\operatorname{Spec}(Z(Q_1))\times_S\operatorname{Spec}(Z(Q_2))\longrightarrow\operatorname{Spec}(Z(Q)),
\]
使得在任意\Term{base-change}{基变换}之后，都有
\[
e(s+t)=e(s)e(t)+(1-e(s))(1-e(t)).
\]
若 $s(W)$ 是满足 $e(s(W))=e(W)$ 的 $\operatorname{Spec}(Z(Q))$ 的\Term{section}{截面}，则公式 1.10.1 可改写为
\[
s(W_1\oplus W_2)=s(W_1)+s(W_2).
\]

\begin{SGAresult}{命题}{1.12}
设 $A$ 是一个域，$V$ 是 $V$ 上一个\Term{dimension}{维数}为 $2m>0$ 的\Term{vector-space}{向量空间}，$Q$ 是 $V$ 上一个\Term{nondegenerate}{非退化}\Term{quadratic-form}{二次型}，$W_1,W_2$ 是 $V$ 的两个\Term{dimension}{维数}为 $m$ 的\Term{totally-isotropic}{全迷向}\Term{vector-subspace}{子空间}。则 $e(W_1)=e(W_2)$ 当且仅当 $\dim(W_1/(W_1\cap W_2))$ 为偶数。
\end{SGAresult}
\SourcePage{80}
\begin{proof}
（参见 Bourbaki Alg 9，§6，习题 28。）设 $e_1\cdots e_m f_1\cdots f_m$ 是 $V$ 的一个基，使得

\noindent a) $Q(\sum x_ie_i+\sum y_if_i)=\sum x_iy_i$；

\noindent b) $W_1$ 由各个 $e_i$ 生成；

\noindent c) $W_2$ 由 $(e_i)_{1\le i\le s}$ 及 $(f_i)_{s<i\le m}$ 生成。

于是 $V$ 是各个 $V_i=Ae_i+Af_i$ 的和，而且
\[
W_\alpha=\bigoplus W_\alpha\cap V_i\qquad(\alpha=1,2).
\]
因而加法公式 1.11 可将问题归结到 $\dim V=2$ 的情形。若 $W_1=W_2$，断言显然。反之，若 $W_1=Ae$、$W_2=Af$，且 $\Phi(e,f)=1$，则
\[
e(W_1)=fe,\qquad e(W_2)=ef=1-fe.
\]
\end{proof}

\SGAsection{2}{二次超曲面}
\textbf{2.1。} 设 $k$ 是\Term{algebraically-closed}{代数闭}域，$V$ 是 $k$ 上\Term{dimension}{维数}为 $n+2$ 的\Term{vector-space}{向量空间}。$k$ 上一个\Term{dimension}{维数}为 $n$ 的\Term{smooth}{光滑}\Term{quadric}{二次超曲面}，是指 $k$ 上一个\Term{scheme}{概形} $X$，它\Term{isomorphism}{同构}于 $\mathbb P(V^*)$ 中由方程 $Q=0$ 定义的\Term{subscheme}{子概形}，其中 $Q$ 是 $k$ 上的一个\Term{ordinary-quadratic-form}{普通二次型}。

\begin{SGAresult}{注记}{2.2}
对 $n=0,1,2$，$X$ 是 $k$ 上\Term{dimension}{维数}为 $n$ 的\Term{smooth}{光滑}\Term{quadric}{二次超曲面}，当且仅当分别有
\[
\begin{array}{ll}
n=0:&X\simeq\operatorname{Spec}(k)\amalg\operatorname{Spec}(k),\\
n=1:&X\simeq\mathbb P^1_k,\\
n=2:&X\simeq\mathbb P^1_k\times\mathbb P^1_k.
\end{array}
\]
\end{SGAresult}

\SourcePage{81}
\begin{SGAresult}{\Term{lemma}{引理}}{2.3}
沿用 2.1 的记号，设 $X$ 由 $Q=0$ 定义。

\noindent (i) $\Omega_X^n\simeq\mathscr O_X(-n)$。

\noindent (ii)
\[
\operatorname{Pic}(X)\simeq
\begin{cases}
0,&n=0,\\
\mathbb Z,&n=1;\quad\mathscr O_X(1)\text{ 是一个生成元的平方},\\
\mathbb Z\oplus\mathbb Z,&n=2,\\
\mathbb Z,&n\ge3;\quad\mathscr O_X(1)\text{ 是一个生成元}.
\end{cases}
\]
\noindent (iii) 对 $i>0$ 有 $H^i(X,\mathscr O_X)=0$，且 $H^i(X,\mathscr O_X(1))=0$。

\noindent (iv)
\[
H^0(X,\mathscr O_X(1))\xleftarrow{\sim}H^0(\mathbb P(V^*),\mathscr O_X(1))=V^*,
\]
并且当 $n>0$ 时，$H^0(X,\mathscr O_X)=k$。
\end{SGAresult}
\begin{proof}
断言 (i) 来自以下事实：$X$ 在 $\mathbb P(V^*)$ 中的\Term{conormal-sheaf}{余法层}是 $\mathscr O_X(-2)$，而 $\mathbb P^r$ 的\Term{canonical-sheaf}{典范层}是 $\mathscr O(-r-1)$。当 $n\le2$ 时，断言 (ii) 来自 2.2；当 $n\ge3$ 时，它来自 SGA 2 第 XII 讲 3.7。最后，(iii) 和 (iv) 来自 FAC n\textsuperscript{o}78。
\end{proof}

\begin{SGAresult}{定义}{2.4}
\Term{scheme}{概形} $S$ 上一个\Term{dimension}{维数}为 $n$ 的\Term{smooth}{光滑}\Term{quadric}{二次超曲面}，是指在 $S$ 上\Term{proper}{固有}且\Term{smooth}{光滑}、\Term{geometric-fibre}{几何纤维}为\Term{smooth}{光滑}\Term{quadric}{二次超曲面}的 $S$-\Term{scheme}{概形} $f:X\to S$。
\end{SGAresult}
当 $n=0$ 时，$S$ 上\Term{dimension}{维数}为 $0$ 的\Term{smooth}{光滑}\Term{quadric}{二次超曲面}就是 $S$ 的一个双重\Term{etale-covering}{平展覆盖}。当 $n=1$ 时，它就是 $S$ 上一个\Term{dimension}{维数}为 $1$ 的\Term{severi-brauer-variety}{Severi--Brauer 簇}。

\textbf{2.5。} 若 $p:X\to S$ 是 $S$ 上一个\Term{dimension}{维数}为 $n$ 的\Term{smooth}{光滑}\Term{quadric}{二次超曲面}，则根据 2.3 (i)，\Term{invertible-sheaf}{可逆层} $\Omega_{X/S}^n$ 的逆是\Term{ample}{丰沛}的。

\SourcePage{82}
根据 2.3 (ii)，\Term{picard-scheme}{Picard 概形} $\operatorname{Pic}_{X/S}$ 在 $S$ 上\Term{smooth}{光滑}（因为 $H^2(X,\mathscr O_X)=0$），在 $S$ 上\Term{unramified}{非分歧}（因为 $H^1(X,\mathscr O_X)=0$），并且在 $S$ 上本质\Term{proper}{固有}（因为 $X/S$ \Term{smooth}{光滑}）。因此，它在 $S$ 上按\Term{etale-topology}{平展拓扑}\Term{local-scheme}{局部}地是一个\Term{discrete}{离散}\Term{constant-group}{常值群}，其结构如 2.3 (iii) 所示。

特别地，若 $X$ 有一个\Term{section}{截面}，则 $X$ 上存在一个\Term{invertible-sheaf}{可逆层} $L$，使 $L^{\otimes(-n)}$ 与 $\Omega_{X/S}^n$ 在 $\operatorname{Pic}_{X/S}$ 中有相同的像；任意两个这样的\Term{sheaf}{层}在 $S$ 上\Term{local-scheme}{局部}\Term{isomorphism}{同构}。根据 2.3 (iii) 及 [1] 0.5，第19页，$p_*L$ \Term{locally-free}{局部自由}，且其形成与任意\Term{base-change}{基变换}相容。因此 $p_*L$ 是\Term{rank}{秩}为 $n+2$ 的\Term{locally-free}{局部自由}模（2.3 (iv)），从 $p^*p_*L$ 到 $L$ 的映射\Term{surjectivity}{满射}，从 $X$ 到 $\mathbb P(p_*L)$ 的\Term{canonical-map}{典范映射}为\Term{closed-immersion}{闭浸入}；它将 $X$ 与 $\mathbb P(p_*L)$ 中一个\Term{degree}{次数}为 $2$ 的\Term{relative-divisor}{相对除子}等同，并且在 $S$ 上\Term{local-scheme}{局部}地，$X$ 是由 $(p_*L)^*$ 上一个\Term{ordinary-quadratic-form}{普通二次型}定义的 $\mathbb P(p_*L)$ 中的\Term{quadric}{二次超曲面}；这个\Term{quadratic-form}{二次型}在相差一个\Term{invertible}{可逆}因子的意义下唯一。

\textbf{2.6。} 假设 $n>0$。若 $L_1,L_2$ 是上述类型的两个\Term{invertible-sheaf}{可逆层}，则在 $S$ 上\Term{local-scheme}{局部}地存在一个\Term{isomorphism}{同构} $u:L_1\xrightarrow{\sim}L_2$。由于 $p_*\mathscr O_X=\mathscr O_S$（根据 2.3 (iv)），这一\Term{isomorphism}{同构}在相差一个\Term{invertible}{可逆}因子的意义下唯一。因此，各个\Term{isomorphism}{同构} $u:L_1\xrightarrow{\sim}L_2$ 诱导同一个\Term{isomorphism}{同构}
\[
\mathbb P(u):\mathbb P(p_*L_1)\xrightarrow{\sim}\mathbb P(p_*L_2).
\]
所以 $\mathbb P(p_*L_1)$ \Term{canonical}{典范}\Term{isomorphism}{同构}于 $\mathbb P(p_*L_2)$，且仅依赖于 $X/S$，与 $L$ 的选择无关。将它记作 $P(X)$。

\SourcePage{83}
\Term{projective-space}{射影空间} $P(X)$ 在 $S$ 上按\Term{etale-topology}{平展拓扑}\Term{local-scheme}{局部}地定义。由\Term{descent}{下降}，它定义了 $S$ 上一个\Term{severi-brauer-variety}{Severi--Brauer 簇}，仍记作 $P(X)$；$X$ 是其中一个\Term{degree}{次数}为 $2$ 的\Term{relative-divisor}{相对除子}。

在 $S$ 上\Term{local-scheme}{局部}地，这对\Term{object}{对象} $X\hookrightarrow P(X)$ 定义了一个\Term{clifford-algebra}{Clifford 代数} $C^+(X,P(X))$（1.3）。由\Term{descent}{下降}，得到 $S$ 上一个\Term{algebra}{代数}，简记为 $C^+(X)$。

\textbf{2.7。} 假设 $n=2m>0$。于是 $C^+(X)$ 的\Term{centre-algebra}{中心} $Z$ 是一个 $\mathscr O_S$-\Term{algebra}{代数}，定义 $S$ 的一个双重\Term{etale-covering}{平展覆盖} $Z(X)$。

$X$ 的一个\Term{quadric-generator}{母空间}（génératrice），是指 $X$ 中一个\Term{subscheme}{子概形} $D$，它是 $P(X)$ 中一个\Term{dimension}{维数}为 $m$ 的\Term{linear-subspace}{线性子空间}。

在 $S$ 上\Term{local-scheme}{局部}地，当 $P(X)\simeq\mathbb P(V^*)$ 且 $X$ 由方程 $Q=0$ 定义时，按照定义，$X$ 的一个\Term{quadric-generator}{母空间} $D$ 可视为 $V$ 中一个\Term{local-scheme}{局部}为\Term{direct-summand}{直和因子}、\Term{totally-singular}{全奇异}且\Term{rank}{秩}为 $m+1$ 的\Term{submodule}{子模} $W(D)$。根据 1.8，$C^+(X)$ 的\Term{centre-algebra}{中心}中的\Term{idempotent}{幂等元} $e(W(D))$ 仅依赖于 $D\subset X$。以 $e(D)$ 表示 $Z(X)$ 的相应\Term{section}{截面}（即 $e(W(D))$ 在其上等于 $1$ 的\Term{section}{截面}）。

$X/S$ 的\Term{quadric-generator}{母空间}是一个在 $S$ 上\Term{projective-morphism}{射影}的\Term{scheme}{概形} $\operatorname{Gén}(X)$ 的 $S$-点；$e$ 定义一个 $S$-\Term{scheme}{概形}\Term{morphism}{态射}
\[
e:\operatorname{Gén}(X)\longrightarrow Z(X).\tag{2.7.1}
\]
当 $n=0$ 时，置 $Z(X)=X$，称 $X$ 的一个\Term{section}{截面}为 $X$ 的一个\Term{quadric-generator}{母空间}，并定义 $e$ 为\Term{identity-map}{恒等映射}。

\SourcePage{84}
\begin{SGAresult}{命题}{2.8}
在 2.7 的假设下，$p:\operatorname{Gén}(X)\to S$ \Term{projective-morphism}{射影}且\Term{smooth}{光滑}，而 (2.7.1) 是它的\Term{stein-factorization}{Stein 分解}。
\end{SGAresult}
\begin{proof}
问题对 $S$ 上的\Term{etale-topology}{平展拓扑}是\Term{local-scheme}{局部}的。\Term{local-scheme}{局部}地，若 $D$ 是一个\Term{quadric-generator}{母空间}，则根据 1.6，与 $D$ 不交的\Term{quadric-generator}{母空间} $D'$ 可与 $S$ 上某个适当\Term{affine-space}{仿射空间}的 $S$-点等同；由此可用在 $S$ 上\Term{smooth}{光滑}的\Term{open-subset}{开集}\Term{covering}{覆盖} $\operatorname{Gén}(X)$。

为证明 $e$ 的\Term{geometric-fibre}{几何纤维}是\Term{geometrically-connected}{几何连通}的，只须考虑 $S$ 为一个\Term{algebraically-closed}{代数闭}域 $k$ 的\Term{spectrum}{谱}的情形，其中 $X$ 由 $\mathbb P(V^*)$ 中的方程 $Q=0$ 定义。

因此设 $W_1$ 和 $W_2^*$ 是 $V$ 中两个\Term{rank}{秩}为 $m+1$ 的\Term{totally-isotropic}{全迷向}\Term{vector-subspace}{子空间}，满足 $e(W)=e(W')$，即 $W/(W\cap W')$ 的\Term{dimension}{维数}为偶数（1.12）。像 1.12 中那样，可找到 $V$ 的一个\Term{orthogonal-decomposition}{正交分解}
\[
V=V_0\oplus\bigoplus V_i,
\]
使得
\[
\begin{array}{ll}
\text{a)}&W_\alpha=W_\alpha\cap V_0\oplus\bigoplus W_\alpha\cap V_i\quad(\alpha=1,2),\\
\text{b)}&W_1\cap V_0=W_2\cap V_0,\\
\text{c)}&\dim V_i=4,\quad V_i=(W_1\cap V_i)\oplus(W_2\cap V_i).
\end{array}
\]
只须证明 $W_1\cap V_i$ 和 $W_2\cap V_i$ 属于 $V_i$ 中极大\Term{totally-singular}{全奇异}\Term{vector-subspace}{子空间}的同一个\Term{connected}{连通}族。换言之，归结到了\Term{dimension}{维数}为 $2$ 的\Term{quadric}{二次超曲面}的情形。
\SourcePage{85}
这样的\Term{quadric}{二次超曲面}\Term{isomorphism}{同构}于 $\mathbb P^1\times\mathbb P^1$，其\Term{quadric-generator}{母空间}为\Term{subscheme}{子概形} $\{t\}\times\mathbb P^1$ 和 $\mathbb P^1\times\{t\}$；此时命题显然。
\end{proof}

\SGAsection{3}{二次超曲面的上同调}
\textbf{3.1。} 设 $p:X\to S$ 是 $S$ 上一个\Term{dimension}{维数}为 $n$ 的\Term{smooth}{光滑}\Term{quadric}{二次超曲面}，$\ell$ 是在 $S$ 上\Term{invertible}{可逆}的素数。以
\[
\eta\in H^0(S,R^2p_*\mathbb Z_\ell(1))
\]
表示“\Term{hyperplane-section}{超平面截面}的\Term{cohomology-class}{上同调类}”。在 $S$ 上按\Term{etale-topology}{平展拓扑}\Term{local-scheme}{局部}地，有 $P(X)\simeq\mathbb P(V^*)$；此时按照定义，$\eta$ 是\Term{first-chern-class}{第一 Chern 类}
\[
c_1(\mathscr O(1))\in H^2(X,\mathbb Z_\ell(1))
\]
在 $H^0(S,R^2p_*\mathbb Z_\ell(1))$ 中的像。

\textbf{3.2。} 假设 $n=2m$。于是 $X$ 的一个\Term{quadric-generator}{母空间} $D$ 定义了 $H^n(X,\mathbb Z_\ell(m))$ 中的一个\Term{cohomology-class}{上同调类} $\operatorname{cl}'(D)$，其在 $H^0(S,R^np_*\mathbb Z_\ell(m))$ 中的像记为 $\operatorname{cl}(D)$。由 2.8 可知，后者仅依赖于 $Z(X)$ 的\Term{section}{截面} $e(D)$，从而得到 $S$ 上\Term{etale-sheaf}{平展层}的一个\Term{morphism}{态射}
\[
\operatorname{cl}:Z(X)\longrightarrow R^np_*\mathbb Z_\ell(m).\tag{3.2.1}
\]

\begin{SGAresult}{\Term{theorem}{定理}}{3.3}
设 $p:X\to S$ 是 $S$ 上一个\Term{dimension}{维数}为 $n$ 的\Term{smooth}{光滑}\Term{quadric}{二次超曲面}，$\ell$ 是在 $S$ 上\Term{invertible}{可逆}的素数。

\noindent (i) $R^{2i+1}p_*\mathbb Z_\ell=0$。
\SourcePage{86}

\noindent (ii) 若 $0\le2i<n$（相应地，若 $n<2i\le2n$），则 $R^{2i}p_*\mathbb Z_\ell(i)$ \Term{canonical}{典范}\Term{isomorphism}{同构}于\Term{constant-sheaf}{常值层} $\mathbb Z_\ell$，且由 $\eta^i$（相应地，由 $\eta^i/2$）生成。

\noindent (iii) 若 $n=2m$，则由 (3.2.1) 得到的映射
\[
\mathbb Z_\ell^{Z(X)}\longrightarrow R^np_*\mathbb Z_\ell(m)\tag{3.2.1}
\]
是\Term{isomorphism}{同构}；若 $\alpha,\beta$ 是 $Z(X)$ 的两个不交\Term{section}{截面}（假设 $Z(X)\simeq S\amalg S$；\Term{local-scheme}{局部}地总是如此），则有
\[
\begin{aligned}
\text{(a)}\quad&\eta^m=\operatorname{cl}(\alpha)+\operatorname{cl}(\beta);\\
\text{(b)}\quad&m\text{ 为偶数时：}\quad
\operatorname{Tr}(\operatorname{cl}(\alpha)^2)=\operatorname{Tr}(\operatorname{cl}(\beta)^2)=1,
\quad\operatorname{cl}(\alpha)\operatorname{cl}(\beta)=0;\\
&m\text{ 为奇数时：}\quad\operatorname{cl}(\alpha)^2=\operatorname{cl}(\beta)^2=0,
\quad\operatorname{Tr}(\operatorname{cl}(\alpha)\operatorname{cl}(\beta))=1;\\
\text{(c)}\quad&\eta\,(\operatorname{cl}(\alpha)-\operatorname{cl}(\beta))=0.
\end{aligned}
\]
\end{SGAresult}
\begin{proof}
断言 (i)、(ii) 仅为便于查考而列出（第 XI 讲 1.6；对于 (i) 中 $2i+1=n$ 的情形，还须用第 XI 讲 2.6 补充）。

只须在 $S$ 是一个\Term{algebraically-closed}{代数闭}域 $k$ 的\Term{spectrum}{谱}时验证 (iii)（甚至只取 $S=\operatorname{Spec}(\mathbb C)$ 也足够）。

\textbf{公式 (a)。} 可以假设 $X$ 是 $\mathbb P^{m+1}(k)$ 中方程 $\sum_0^m x_ix_{i+m+1}=0$ 的\Term{quadric}{二次超曲面}。类 $\eta^m$ 是方程
\[
x_i=0\qquad(0\le i<m)
\]
所定义的平面\Term{section}{截面}的\Term{cohomology-class}{上同调类}，而这一\Term{cycle}{圈}是下列方程定义的两个\Term{quadric-generator}{母空间} $D_1,D_2$ 的和：
\SourcePage{87}
\[
\begin{array}{ll}
D_1:&x_i=0\quad(0\le i\le m),\\
D_2:&x_i=0\quad(0\le i<m),\quad x_{2m+1}=0.
\end{array}
\]
根据 1.12，$e(D_1)\ne e(D_2)$，所以
\[
\eta^m=\operatorname{cl}(D_1)+\operatorname{cl}(D_2)=\operatorname{cl}(\alpha)+\operatorname{cl}(\beta).
\]

\textbf{公式 (b)。} 若\Term{quadric-generator}{母空间} $D_1,D_2$ 不交（相应地，在一点\Term{transverse-intersection}{横截相交}），则
\[
\operatorname{cl}(D_1)\operatorname{cl}(D_2)=0
\quad\text{（相应地，}\operatorname{Tr}(\operatorname{cl}(D_1)\operatorname{cl}(D_2))=1\text{）。}
\]
根据 1.12，当 $m$ 为偶数时，有 $e(D_1)\ne e(D_2)$（相应地，$e(D_1)=e(D_2)$）；当 $m$ 为奇数时，有 $e(D_1)=e(D_2)$（相应地，$e(D_1)\ne e(D_2)$）。公式 (b) 随即得到。

\textbf{公式 (c)。} 当 $n=0$ 时显然。否则，根据 (ii)，只须证明 $\eta^m(\operatorname{cl}(\alpha)-\operatorname{cl}(\beta))=0$。事实上，
\[
\operatorname{Tr}\bigl(\eta^m(\operatorname{cl}(\alpha)-\operatorname{cl}(\beta))\bigr)
=\operatorname{Tr}\bigl(\operatorname{cl}(\alpha)^2-\operatorname{cl}(\beta)^2\bigr)=0.
\]

\textbf{核对。}
\[
\begin{aligned}
m\text{ 为偶数：}\quad
2&=\operatorname{tr}(\eta^2)\\
&=\operatorname{tr}\bigl(\operatorname{cl}(\alpha^2)+2\operatorname{cl}(\alpha)\operatorname{cl}(\beta)+\operatorname{cl}(\beta)^2\bigr)=1+1;\\
m\text{ 为奇数：}\quad
2&=\operatorname{tr}(\eta^2)\\
&=\operatorname{Tr}\bigl(\operatorname{cl}(\alpha)^2+2\operatorname{cl}(\alpha)\operatorname{cl}(\beta)+\operatorname{cl}(\beta)^2\bigr)=2.
\end{aligned}
\]
现在，3.3.1 是\Term{isomorphism}{同构}这一事实，由下述二变量\Term{bilinear-form}{双线性型}
\[
B(X,Y)=x_1y_2+x_2y_1\quad\text{和}\quad B(x,y)=x_1y_1+x_2y_2
\]
的\Term{discriminant}{判别式}为 $\pm1$ \Term{pushout}{推出}。
\end{proof}

\SourcePage{88}
\begin{SGAresult}{核对}{3.4}
设 $X$ 是\Term{finite-field}{有限域} $\mathbb F_q$ 上一个\Term{dimension}{维数}为 $n$ 的\Term{smooth}{光滑}\Term{quadric}{二次超曲面}。结合 3.3，\Term{lefschetz-trace-formula}{Lefschetz 迹公式}给出经典公式
\[
\#X(\mathbb F_q)=
\begin{cases}
\displaystyle\sum_{0}^{n}q^i,&n\text{ 为奇数},\\
\displaystyle\sum_{0}^{n}q^i+\varepsilon q^m,&n=2m\quad(\varepsilon=\pm1).
\end{cases}
\]
当 $n$ 为偶数时，若 $X$ \Term{split-quadric}{分裂}，即有一个有理\Term{quadric-generator}{母空间}，则 $\varepsilon=1$；否则 $\varepsilon=-1$。
\end{SGAresult}

\textbf{3.5。} 若 $X$ 是 $S$ 上一个\Term{dimension}{维数}为 $n=2m$ 的\Term{smooth}{光滑}\Term{quadric}{二次超曲面}，则称 $R^nf_*\mathbb Z_\ell(m)$ 中 $\eta^m$ 的\Term{orthogonality}{正交}\Term{space}{空间}为它的\Term{primitive-part}{本原部分}（沿用 3.3 的记号，此\Term{orthogonality}{正交}\Term{space}{空间}由 $\operatorname{cl}(\alpha)-\operatorname{cl}(\beta)$ 生成）。称 $R^nf_*\mathbb Z_\ell(m)$ 模掉 $\mathbb Z_\ell\eta^m$ 的商为 $R^nf_*\mathbb Z_\ell(m)$ 的\Term{primitive-quotient}{本原商}。将这些\Term{object}{对象}与 $\mathbb Z_\ell(-m)$ 作\Term{tensor-product}{张量积}，即定义了 $R^nf_*\mathbb Z_\ell$ 的\Term{primitive-part}{本原部分}及\Term{primitive-quotient}{本原商}。

\textbf{3.6。} 设 $V$ 是 $S$ 上一个\Term{rank}{秩}为 $n+2$ 的\Term{locally-free}{局部自由}\Term{vector-bundle}{向量丛}，$X$ 是 $V$ 上一个\Term{ordinary-quadratic-form}{普通二次型} $Q$ 所定义的、$S$ 上的\Term{smooth}{光滑}\Term{quadric}{二次超曲面}。再设 $H$ 是 $\mathbb P(V^*)$ 中一个与 $X$ \Term{transverse-intersection}{横截相交}的\Term{hyperplane}{超平面}，并置 $X^0=X-H$。\Term{scheme}{概形} $X$ 及 $Y=X\cap H$ 是 $S$ 上的\Term{smooth}{光滑}\Term{quadric}{二次超曲面}，\Term{dimension}{维数}分别为 $n$ 及 $n-1$。这使我们可以用 3.3 计算 $X^0=X-Y$ 的\Term{cohomology}{上同调}。

\SourcePage{89}
\[
\begin{tikzcd}
X^0\arrow[r,hook]\arrow[dr,"f"']&X\arrow[d,"p"]&Y\arrow[l,hook']\arrow[dl,"q"]\\
&S&
\end{tikzcd}\tag{3.6.1}
\]
若 $n=0$，则 $Y=\varnothing$ 且 $X=X^0$。假设 $n>0$。有两个彼此\Term{dual}{对偶}的\Term{long-exact-sequence}{长正合列}
\[
\cdots\xrightarrow{\partial}R^if_!\mathbb Z_\ell\longrightarrow R^ip_*\mathbb Z_\ell
\xrightarrow{r_i}R^iq_*\mathbb Z_\ell\longrightarrow\cdots,\tag{3.6.2}
\]
\[
\cdots\xrightarrow{\partial}R^{i-2}q_*\mathbb Z_\ell(-1)\longrightarrow R^ip_*\mathbb Z_\ell
\longrightarrow R^if_*\mathbb Z_\ell\longrightarrow\cdots.\tag{3.6.3}
\]
当 $n$ 为偶数时，\Term{restriction}{限制}\Term{homomorphism}{同态} $r_i$ 对 $i\ne n,2n$ 是\Term{isomorphism}{同构}；当 $n$ 为奇数时，它对 $i\ne n-1$ 是\Term{isomorphism}{同构}。对 $n=2m$，若 $\alpha$ 是一个\Term{quadric-generator}{母空间}的\Term{cohomology-class}{上同调类}，则 $r_n(\alpha)=\tfrac12\eta^m$，因此 $r_n$ \Term{surjectivity}{满射}，其\Term{kernel}{核}为 $R^np_*\mathbb Z_\ell$ 的\Term{primitive-part}{本原部分}。对 $r=2m+1$，有 $r_{2m}(\eta^m)=\eta^m$，所以 $r_{2m}$ \Term{injectivity}{单射}，其\Term{cokernel}{余核}为 $R^{2m}q_*\mathbb Z_\ell$ 的\Term{primitive-quotient}{本原商}。由此\Term{pushout}{推出}，当 $i\ne n,2n$ 时，$R^if_!\mathbb Z_\ell=0$；当 $i=n,2n$ 时，$R^if_!\mathbb Z_\ell$ 是一个\Term{rank}{秩}为 $1$ 的自由\Term{locally-constant}{局部常值} $\mathbb Z_\ell$-\Term{sheaf}{层}。由\Term{duality}{对偶性}，或借助 3.5.3，同样可以证明：当 $i\ne0,n$ 时，$R^if_*\mathbb Z_\ell=0$；$f_*\mathbb Z_\ell=\mathbb Z_\ell$；且 $R^nf_*\mathbb Z_\ell$ 的\Term{rank}{秩}为一。

假设 $n=2m$，并沿用 3.3 的记号。
\SourcePage{90}
\[
\begin{tikzcd}[column sep=small]
0\arrow[r]&R^nf_!\mathbb Z_\ell(m)\arrow[r]\arrow[d,"\varphi"']&R^np_*\mathbb Z_\ell(m)\arrow[r]\arrow[d,equal]&R^nq_*\mathbb Z_\ell(m)\arrow[r]&0\\
0&R^nf_*\mathbb Z_\ell(m)\arrow[l]&R^np_*\mathbb Z_\ell(m)\arrow[l]&R^{n-2}q_*\mathbb Z_\ell(m-1)\arrow[l]&0\arrow[l]
\end{tikzcd}
\]
于是\Term{sheaf}{层} $R^nf_!\mathbb Z_\ell(m)$ 有两个互为相反数的自然\Term{generator}{生成元} $\pm\delta$，其在 $R^np_*\mathbb Z_\ell(m)$ 中的像为 $\pm(\operatorname{cl}(\alpha)-\operatorname{cl}(\beta))$。有（3.3 (iii) b）
\[
\operatorname{Tr}(\delta^2)=(\operatorname{cl}(\alpha)-\operatorname{cl}(\beta))^2=(-1)^m\cdot2,
\]
并且 $\pm\varphi(\delta)$ 是 $R^nf_*\mathbb Z_\ell(m)$ 中自然的\Term{dual}{对偶}\Term{generator}{生成元} $\pm\delta'$ 的两倍；这些\Term{generator}{生成元}就是 $\operatorname{cl}(\alpha)$ 和 $\operatorname{cl}(\beta)$ 的（互为相反数的）像。

假设 $n=2m+1$，并对 $Y$ 使用 3.3 的记号。
\[
\begin{tikzcd}[column sep=small]
0\arrow[r]&R^{2m}p_*\mathbb Z_\ell(m)\arrow[r]&R^{2m}q_*\mathbb Z_\ell(m)\arrow[r,"\partial"]&R^nf_!\mathbb Z_\ell(m)\arrow[r]\arrow[d,"\varphi"]&0\\
&R^{2m+2}p_*\mathbb Z_\ell(m)&R^{2m}q_*\mathbb Z_\ell(m-1)\arrow[l]&R^nf_*\mathbb Z_\ell(m)\arrow[l]&0\arrow[l]
\end{tikzcd}
\]
于是\Term{sheaf}{层} $R^nf_!\mathbb Z_\ell(m)$ 有两个互为相反数的自然\Term{generator}{生成元} $\pm\delta$：$\partial\operatorname{cl}(\alpha)$ 及 $\partial\operatorname{cl}(\beta)$。有
\[
\delta\,\varphi(\delta)=\delta^2=(\partial\operatorname{cl}(\alpha))^2
=\partial\bigl(\operatorname{cl}(\alpha)\,\partial\operatorname{cl}(\alpha)\bigr)=0,
\]
所以 $\varphi(\delta)=0$。另一方面，$R^nf_*\mathbb Z_\ell(m+1)$ 有两个互为相反数的自然\Term{generator}{生成元} $\pm\delta'$，其在 $R^{2m}q_*\mathbb Z_\ell(m)$ 中的像为 $\pm(\operatorname{cl}(\alpha)-\operatorname{cl}(\beta))$，并且
\[
\operatorname{Tr}(\delta\delta')=\pm1.
\]

\SourcePage{91}
这些结果汇集在下表中。

\begin{SGAresult}{表}{3.7}
\Term{dimension}{维数}为 $n$ 的\Term{affine-quadric}{仿射二次超曲面} $X^0=X-Y$ 的\Term{cohomology}{上同调}。

\noindent\textbf{A。} \Term{cohomology}{上同调}\Term{torsion-free}{无挠}。非零的\Term{betti-number}{Betti 数}为：
\[
\begin{array}{ll}
\text{无支集限制的上同调：}&b_0=b_n=1,\\
\text{具有固有支集的上同调：}&b_{2n}=b_n=1,
\end{array}
\]
但 $n=0$ 时例外：$b_0=2$。

\noindent\textbf{B。}
\[
\begin{array}{lll}
n=2m>0:&R^nf_!\mathbb Z_\ell(m)&=R^{2m}p_*\mathbb Z_\ell(m)\text{ 的本原部分},\\
&R^nf_*\mathbb Z_\ell(m)&=R^{2m}p_*\mathbb Z_\ell(m)\text{ 的本原商};\\[3pt]
n=2m+1:&R^nf_!\mathbb Z_\ell(m)&=R^{2m}q_*\mathbb Z_\ell(m)\text{ 的本原商},\\
&R^nf_*\mathbb Z_\ell(m+1)&=R^{2m}q_*\mathbb Z_\ell(m)\text{ 的本原部分}.
\end{array}
\]
这些群\Term{local-scheme}{局部}地具有在相差一个\Term{signature}{符号}的意义下确定的自然\Term{generator}{生成元}；对\Term{proper-support-cohomology}{具有固有支集的上同调}群，记作 $\delta$，对其他群记作 $\delta'$。

\noindent\textbf{C。} $\varphi:R^nf_!\mathbb Z_\ell\to R^nf_*\mathbb Z_\ell$ 在 $n$ 为奇数时为零；在 $n>0$ 为偶数时，将 $\pm\delta$ 映为 $\pm2\delta'$。有
\[
\operatorname{Tr}(\delta\delta')=\pm1,
\]
\[
\operatorname{Tr}(\delta^2)=
\begin{cases}
0,&n=2m+1,\\
(-1)^m\cdot2,&n=2m.
\end{cases}
\]
\end{SGAresult}

\begin{SGAresult}{核对}{3.8}
$\mathbb C^n$ 中方程 $\sum z_i^2=1$ 所定义的复\Term{affine-quadric}{仿射二次超曲面} $S$，\Term{diffeomorphism}{微分同胚}于实\Term{sphere}{球面} $S\cap\mathbb R^n$ 的\Term{tangent-bundle}{切丛}。这与 A 和 C 一致。
\end{SGAresult}

\section*{参考文献}
\begin{enumerate}
\item D. Mumford，\emph{几何\Term{invariants}{不变量}理论}（\emph{Geometric invariant theory}），Ergebnisse，第34卷，Springer Verlag，1965年。
\end{enumerate}
